%============================= app4.tex================================
\chapter{Damped Least-Squares Inverse Kinematics as Regularized Optimization}
\label{app:dls}

Section~\ref{subsec:sv_dls} states the damped least-squares (DLS) resolved-rate mapping,
Equation~\eqref{eq:sv_dls}, $\dot{\mathbf{q}} = J^\top(JJ^\top+\lambda^2I)^{-1}\mathbf{v}$, as a
known result, citing~\citep{whitney1969resolved,wampler1986manipulator,chiaverini1994review,
nakamura1986inverse}. This appendix derives it from the underlying regularized least-squares
problem and explains its singularity-robustness property, which motivates the manipulability-based
switching of $\lambda$ used alongside it.

%% ============================================================
\section{The Regularized Least-Squares Problem}\label{sec:appD_problem}
%% ============================================================

The undamped resolved-rate problem is to find joint velocities $\dot{\mathbf{q}} \in \mathbb{R}^6$
producing the commanded Cartesian twist $\mathbf{v} \in \mathbb{R}^6$ of Equation~\eqref{eq:sv_twist}
exactly, $J\dot{\mathbf{q}} = \mathbf{v}$, for the manipulator Jacobian $J(\mathbf{q}) \in
\mathbb{R}^{6\times6}$ of Equation~\eqref{eq:sv_jacobian}. Near a kinematic singularity $J$ loses
rank and this equation has no solution, or an ill-conditioned one requiring arbitrarily large joint
speeds for a small Cartesian motion. Damped least squares avoids this by relaxing exact tracking into
a penalized optimization,
\begin{equation}
\dot{\mathbf{q}}^* = \arg\min_{\dot{\mathbf{q}}} \; \lVert J\dot{\mathbf{q}} - \mathbf{v} \rVert^2
+ \lambda^2 \lVert \dot{\mathbf{q}} \rVert^2,
\label{eq:appD_objective}
\end{equation}
trading exact twist tracking (the first term) against a penalty on joint-speed magnitude (the second
term), with $\lambda \geq 0$ setting the trade-off; $\lambda=0$ recovers the undamped problem exactly.

%% ============================================================
\section{Solving the Regularized Problem}\label{sec:appD_solve}
%% ============================================================

Expanding Equation~\eqref{eq:appD_objective} and setting its gradient with respect to
$\dot{\mathbf{q}}$ to zero,
\begin{equation}
\nabla_{\dot{\mathbf{q}}}\left[ \lVert J\dot{\mathbf{q}}-\mathbf{v} \rVert^2 + \lambda^2\lVert
\dot{\mathbf{q}} \rVert^2 \right]
= 2J^\top(J\dot{\mathbf{q}}-\mathbf{v}) + 2\lambda^2\dot{\mathbf{q}} = 0,
\label{eq:appD_gradzero}
\end{equation}
gives the normal equations $(J^\top J + \lambda^2 I)\dot{\mathbf{q}} = J^\top\mathbf{v}$, and, since
$J^\top J + \lambda^2 I$ is positive definite (hence invertible) for any $\lambda > 0$ regardless of
$J$'s rank,
\begin{equation}
\dot{\mathbf{q}}^* = \left( J^\top J + \lambda^2 I \right)^{-1} J^\top \mathbf{v}.
\label{eq:appD_solution1}
\end{equation}
This is a valid, always-invertible solution, but it inverts a matrix in joint space
($6\times6$ here, but $n\times n$ in general for an $n$-joint arm). The form actually used in
Equation~\eqref{eq:sv_dls} inverts in task space instead; the two are shown equal by the identity
\begin{equation}
(J^\top J + \lambda^2 I)\,J^\top = J^\top J J^\top + \lambda^2 J^\top = J^\top (JJ^\top+\lambda^2I),
\label{eq:appD_identity}
\end{equation}
which holds for any matrix $J$ by direct expansion. Left-multiplying Equation~\eqref{eq:appD_identity}
by $(J^\top J+\lambda^2I)^{-1}$ and right-multiplying by $(JJ^\top+\lambda^2I)^{-1}$ gives
\begin{equation}
(J^\top J+\lambda^2I)^{-1}J^\top = J^\top(JJ^\top+\lambda^2I)^{-1},
\label{eq:appD_pushthrough}
\end{equation}
so substituting the right-hand side of Equation~\eqref{eq:appD_pushthrough} into
Equation~\eqref{eq:appD_solution1} gives exactly Equation~\eqref{eq:sv_dls},
\begin{equation}
\dot{\mathbf{q}}^* = J^\top\left( JJ^\top+\lambda^2I \right)^{-1}\mathbf{v}.
\label{eq:appD_solution2}
\end{equation}
Equations~\eqref{eq:appD_solution1} and~\eqref{eq:appD_solution2} are algebraically identical, not
merely approximately so; Equation~\eqref{eq:appD_solution2} is preferred in implementation because it
inverts a $6\times6$ task-space matrix rather than an $n\times n$ joint-space matrix, which is
cheaper whenever the manipulator has more joints than task-space degrees of freedom.

%% ============================================================
\section{Singularity Robustness and the Manipulability-Based Switch}\label{sec:appD_singularity}
%% ============================================================

Near a singularity, $J$ loses rank and $JJ^\top$ (equivalently $J^\top J$) becomes singular or
near-singular, so the undamped solution $\dot{\mathbf{q}} = J^\top(JJ^\top)^{-1}\mathbf{v}$ blows up
in the direction of the lost degree of freedom. Adding $\lambda^2I$ before inverting keeps
$JJ^\top+\lambda^2I$ bounded away from singularity for any $\lambda>0$: writing $JJ^\top$'s
eigenvalues as $\sigma_i^2 \geq 0$ (the squared singular values of $J$), the corresponding
eigenvalues of $(JJ^\top+\lambda^2I)^{-1}$ are $1/(\sigma_i^2+\lambda^2)$, which stay finite even as
$\sigma_i \to 0$, at the cost of no longer tracking $\mathbf{v}$ exactly along that direction, the
price paid by the regularization term in Equation~\eqref{eq:appD_objective}. This is exactly the
trade-off Section~\ref{subsec:sv_dls} exploits by switching $\lambda$ between a small nominal value
away from singularities and a larger value near them, using the manipulability measure
$\mu(\mathbf{q}) = \sqrt{\det(JJ^\top)} = \prod_i \sigma_i$ as the trigger: $\mu(\mathbf{q}) \to 0$
exactly as the smallest singular value $\sigma_{\min} \to 0$, so $\mu(\mathbf{q})$ is a scalar proxy
for proximity to a singularity without requiring an explicit singular value decomposition at every
control cycle.

%========================================================================

%%% Local  Variables:
%%% mode:  latex
%%% TeX-master: "Thesis.tex"
%%% End:
